% Introduction to the rook endings chapter, for the 1500-1599 book only.
% The theory it needs comes from the library (lib/), placed by make_book where it is first used.

\section*{What to look for}
Rook endings come up more often than any other ending. Their basic theory is short and well known,
and most mistakes in this chapter came from forgetting it. The ideas are explained in the cards
after this introduction; the games show how they were missed.

The defender has these resources:

\begin{itemize}
\item The rook can \term{cut-off}\index{Rook!Cut-off} the enemy king along a rank or file
(\xref{Diagram}{cut-off}). The attacker uses the same idea to keep the defending king away from the
pawn: \xref{Diagram}{kotas-sowa-65w}, \xref{Diagram}{faulhaber-sreyas-65w} and
\xref{Diagram}{meng-zou-56b}.

\item Put the king on the
\term{short-side}\index{Rook!Short side} and the rook on the
\term{long-side},\index{Rook!Long side} on the first or the eighth rank. From there it is ready for
every kind of check.\index{Rook!Good placement of the defender}\uses{good-placement} Many games in this chapter
were lost because the king or the rook went to the wrong side: \xref{Diagram}{kaplan-kara-72b},
\xref{Diagram}{danielsson-haile-76b} and \xref{Diagram}{nespor-kordela-73b}.

\item Check from the right direction:
\term{side-checks}\index{Rook!Side checks} come from the long side (\xref{Diagram}{side-checks},
\xref{Diagram}{meijerink-diegel-66w}),
\term{rear-checks}\index{Rook!Rear checks} from the far end of the
board (\xref{Diagram}{rear-checks}, \xref{Diagram}{hadzala-deng-60b}), and
\term{frontal-checks}\index{Rook!Frontal checks}
stop the pawn from advancing (\xref{Diagram}{faulhaber-sreyas-65w}).

\item With the king in front of the pawn, the
\term{philidor}\index{Rook!Philidor position} holds (\xref{Diagram}{angelov-fil-59w} and
\xref{Diagram}{wasif-adhikari-71b}). A
\term{passive-defence}\index{Rook!Passive defence} with the rook on the first
rank holds only against a knight's or rook's pawn (\xref{Diagram}{alexandersson-mirandaherrera-58w}).

\item Look for a \term{clear-cut-draw}\index{Practical play!Clear-cut draw} first. The simplest one is to
\term{eliminate-the-pawn}:\index{Eliminate the pawn} without it, a rook
ending is a draw at once (\xref{Diagram}{yassibag-tahaoglu-83b}). When nothing else helps, a
\term{desperado}\index{Desperado} can save the game by stalemate (\xref{Diagram}{gloria-blackman-81b}).
\end{itemize}

The attacker has these resources:

\begin{itemize}
\item With the pawn on the seventh rank and the defending king cut off, the attacker wins in
\term{lucena}\index{Rook!Lucena's position} by
\term{bridging}\index{Rook!Bridging}
(\xref{Diagram}{lucena}, \xref{Diagram}{bridging}). See \xref{Diagram}{voges-chen-76w} and
\xref{Diagram}{diazmartin-sanchezramirez-47w}.

\item In the ending of
\term{rook-against-pawn},\index{Rook!Rook against pawn}
the rook wins when its king is close enough or the defending king is cut off from the pawn
(\xref{Diagram}{diazmartin-sanchezramirez-57w} and \xref{Diagram}{uzun-kahya-73w}).

\item Bring the king into play:
\term{king-activity}\index{Strategy!King activity} matters in every ending, and a king left on the
back rank is often why a win slips away (\xref{Diagram}{majstrova-petruzalek-62b} and
\xref{Diagram}{paris-geisler-49w}).

\item When you are winning,
\term{do-not-rush}.\index{Strategy!Do not rush} Several winning positions were spoilt by a pawn push one move
too early (\xref{Diagram}{ilichov-tsintsarski-56w}).
\end{itemize}

Most positions here need no deep calculation. Choose a plan, and rule out what loses with the
\term{process-of-elimination}\index{Practical play!Process of elimination}: see
\xref{Diagram}{vasiliauskas-leboette-63w} for a plan found this way, and
\xref{Diagram}{kahya-incecik-72b} for a
\term{fortress}\index{Fortress} that holds without any
calculation.

